\chapter{The Agentic Paradigm}
\label{chapter:AgenticParadigm}
% EVOLVE-BLOCK-START

An \emph{agent} can be described as a system that perceives its environment, reasons about what to do, and acts repeatedly, over time.\index{Agent}
For most of computing history, this description applied mainly to robots and game-playing programs.
The arrival of LLMs with tool-calling capability changes the picture: a language model that can read files, execute code, browse the web, and call APIs is, in a meaningful sense, an agent that can do knowledge work.
This chapter establishes the vocabulary and conceptual structure the rest of the notes build on.
It opens by defining the agent and its perceive--reason--act loop; Section~\ref{section:CapabilityLadder} then lays out the four-rung \emph{capability ladder} (operate, build, orchestrate, evolve) that orders the course progression; Section~\ref{section:TwoSpines} introduces the two \emph{spines}, evaluation-becomes-fitness and (agentic) context engineering, that thread every rung; and the closing section asks why general-purpose tool-using models make this the moment the pieces combine.

\section{What Is an Agent?}

The core of an agentic system is the \emph{perceive–reason–act loop}.\index{Perceive–reason–act loop}
At each step the agent receives some \emph{observation} from the environment (a file listing, a tool result, a user message, etc.); forms a \emph{plan} or \emph{reasoning trace}; and emits an \emph{action} such as a tool call, a response, a file write.
These three names describe the system's observable behavior, not cognition inside the model: ``reason'' here refers to the intermediate text or plan the loop produces and then acts on, whatever process generated it.
The action changes the environment, producing the next observation.
The loop repeats until the task is complete or the agent determines it cannot proceed.

\begin{definition}
An \emph{agent} is built from four parts, named here because the rest of these notes refer to them: a \emph{model} $M$, which reads the current context and chooses the next action; a set of \emph{tools} $T$ the model can invoke; a \emph{context} $C$, the information the model is given to read at each step; and a \emph{loop} $L$, which runs the model, carries out the action it chooses, feeds the result back as the next observation, and repeats until the task is done or the agent cannot proceed.
\end{definition}

The components interact in a specific way.
The context $C$ is assembled from the system prompt, the task description, a history of previous observations and actions, and any additional information injected by the loop.
The model $M$ reads the full context and selects an action: either a \emph{tool call} (drawn from $T$) or a \emph{final response}.
The loop $L$ executes the tool call, appends the result to the context as a new observation, and invokes $M$ again.

\begin{example}
A concrete trajectory makes the four parts easier to hold onto.
Suppose an ecologist asks an agent whether a published growth model still fits a newly released field dataset.
The tool set $T$ holds three capabilities: list the contents of a directory, read a file, and run a Python script.
The context $C$ begins with the request itself and a note that the data live in a folder called \code{field2026}.
On the first pass the model $M$ reads $C$, has no idea what that folder contains, and calls the directory-listing tool; the loop $L$ runs it and appends the file names as the next observation.
On the second pass $M$ now sees a file named \code{biomass.csv}, calls the file-reading tool on its opening lines, and learns the column names.
On the third it writes a short fitting script, calls the run-Python tool, and reads back the residuals.
Only then does it produce a final answer.
The ecologist never specified this sequence: each step was chosen from what the previous step returned, which is exactly what distinguishes a loop from a recipe.
\end{example}

The model $M$ is not a database and not a classical reasoner; it is a next-token predictor conditioned on $C$, and four properties of that substrate govern where an agent built on it will fail (Anthropic, 2025).
\emph{Next-token prediction} means the model repeatedly estimates a distribution over the next fragment of text given everything in $C$, draws one fragment from it, appends it, and repeats.
There is no stored table of facts it consults and no separate step that checks a claim before writing it down.
It generates by statistical continuation, so it produces fluent text that may be confidently wrong, the failure mode called \emph{hallucination}.
Its knowledge is a frozen snapshot fixed at a training cutoff, blind to anything more recent or too rare to have been well represented in training.
Its working memory is exactly the context window introduced below and developed in Chapter~\ref{chapter:ContextEngineering}: bounded, non-persistent, and unevenly attended.
And its steerability is pattern continuation rather than comprehension of intent: it follows the \emph{form} of an instruction, which is why long chains accumulate error and why text injected through a tool result can hijack a task.
That last failure is worth picturing concretely: a web page an agent has been asked to summarize may itself contain a sentence addressed to the agent, and a system that continues patterns rather than weighing whose instruction it is reading may simply follow it.
Each rung of the ladder that follows is, in part, a discipline for managing one of these four failure modes.

\begin{figure}[htb!]
\begin{center}
%% agent-loop.tex — the perceive-reason-act loop (Chapter 2). House conceptual style.
\input{figures/concept-style}
\begin{tikzpicture}[node distance=1.8cm and 2.2cm]
\node[conceptbox, fill=blue!15]                 (obs)    {Observe};
\node[conceptbox, fill=orange!15, right=of obs]    (reason) {Reason};
\node[conceptbox, fill=green!15,  right=of reason] (act)    {Act};
\draw[conceptflow] (obs)    -- (reason) node[midway, above, conceptlabel] {context};
\draw[conceptflow] (reason) -- (act)    node[midway, above, conceptlabel] {action};
\draw[conceptflow] (act) to[out=270, in=270, looseness=0.6]
      node[above, conceptlabel] {tool result / new state} (obs);
\end{tikzpicture}

\caption{The perceive--reason--act loop.
The agent observes the current state of the environment, reasons to select an action, executes the action via a tool call, and receives the result as the next observation.}
\label{figure:AgentLoop}
\end{center}
\end{figure}

\subsection{Tools}

A \emph{tool}\index{Tool} is any capability the model can invoke as a discrete action: reading a file, running a shell command, calling a web API, querying a database, or invoking another model.
Two items on that list deserve a gloss for readers new to professional tooling.
A \emph{shell command} is an instruction typed at a text prompt that drives the operating system directly, the interface Chapter~\ref{chapter:AgenticRig} takes up in earnest.
An \emph{API}\index{API}, an application programming interface, is a machine-facing entry point to a service: the same archive a website serves to a human, reachable by a program and answering with structured data rather than a rendered page.
A more recent tool type is \emph{computer use}\index{Computer use}: instead of calling a purpose-built API, the model observes a screenshot and issues keyboard and mouse actions, which lets it operate arbitrary desktop software that exposes no programmatic interface.
Tools are described to the model as structured specifications (a name, a natural-language description, and a parameter schema), so the model knows what each tool does and how to call it.
The \emph{parameter schema} is the machine-readable equivalent of a function signature, a declaration of which arguments the tool accepts and of what type: a \code{fetch\_series} tool might declare a station identifier as text and a start and end year as integers, so a malformed request can be caught before anything is executed.
Chapter~\ref{chapter:AgentLoop} treats tool definitions in depth.

The tool set $T$ determines what the agent can \emph{do}, while the model $M$ determines what it \emph{chooses} to do.
A capable model with a sparse tool set is limited by its reach; a rich tool set with a weak model produces unreliable actions.
Designing the right tool set for a task and selecting a model are two of the early decisions in building an agentic environment.

\subsection{Context}

The \emph{context window}\index{Context!window} is the finite-length sequence of tokens the model reads at each invocation.
A \emph{token}\index{Token} is the unit in which the model reads and writes, roughly a common word or a fragment of one; as a rule of thumb English prose runs about three tokens to four words, so a window of two hundred thousand tokens holds a few hundred pages of text.
Everything the agent knows at a given step (its instructions, its history, domain knowledge, tool results) must fit within this window.
Unless managed explicitly, context is not persistent: what the model does not hold in its context window, it does not know.
The practical consequence catches newcomers out: an agent that worked out a calibration constant an hour ago has no memory of it in a fresh session, and will only know it again if something puts it back in front of the model, whether a file it re-reads or a note the loop injects.

This constraint makes the management of context one of the central engineering problems in agentic systems.
Chapter~\ref{chapter:ContextEngineering} develops context engineering as a discipline; Chapter~\ref{chapter:Retrieval} shows how retrieval extends the effective knowledge available within a fixed window.

\begin{remark}
A common misconception is that a larger context window reduces the need for context engineering.
This misconception is as old as large windows themselves; the first hundred-thousand-token windows were pitched as a way to simply feed the model everything, yet in practice models exhibit \emph{position bias} and \emph{context rot} (both developed in Chapter~\ref{chapter:ContextEngineering}).
For instance, information at the beginning and end of the context tends to receive more attention than information buried in the middle.
Context with dissonant content interferes with inference tasks, reasoning, and prediction.
Choosing what to put in the window, and where, remains consequential even as windows grow.
\end{remark}

Not every task needs the full loop.
The distinction between an \emph{agent}, which lets the model dynamically direct its own process and tool use, and a \emph{workflow}, which routes the model through predetermined code paths, is worth drawing early (Anthropic, 2024; Figure~\ref{figure:WorkflowAgentSpectrum}).
A workflow fits predictable, well-scoped tasks whose steps are known in advance; an agent fits open-ended problems where the number of steps cannot be fixed ahead of time, at the price of higher cost and more room for compounding error.
The rungs that follow, and the orchestration patterns of Chapter~\ref{chapter:Orchestration}, move along this spectrum rather than leaping to full autonomy; the durable advice, as of this writing in 2026, is to reach for a workflow when the task allows it and reserve full agency for problems that require it.

\begin{example}
A single laboratory produces both designs in the same week.
Nightly processing of instrument output (load the file, drop the calibration rows, fit the standard curve, write a report) is a workflow: the steps are fixed, their order never varies, and a model may well be used inside one step, say to caption anomalies, without deciding what happens next.
Investigating why last Tuesday's run looks wrong is agentic: nobody can say in advance whether the answer lies in the instrument log, the reagent lot, or a code change, so the sequence of lookups has to be chosen as evidence arrives.
The first belongs in a script that runs unattended; the second is where letting a model steer earns its unpredictability.
\end{example}

These two designs describe the artifact that \emph{runs}; neither describes how it is \emph{built}.
Agentic development, the practice these notes teach, is a build-time activity: a researcher directs an agent such as Claude Code, and the software it produces is frequently a deterministic workflow, or ordinary analysis code, with no agent inside it at all.
Developing \emph{with} an agent is thus distinct from deploying one, and a central aim of these notes is to use the former to build reliable, largely non-agentic research pipelines, reserving in-production agency for the tasks that require it.

\begin{figure}[htb!]
\begin{center}
%% workflow-agent.tex — the workflow/agent distinction (Chapter 2). House conceptual style.
%% Two block nodes on a determinism-to-autonomy spectrum, each with its distinctions listed below.
\input{figures/concept-style}
\begin{tikzpicture}
% the spectrum the two sit on
\draw[conceptflow, <->] (-2.1,1.2) -- (8,1.2);
\node[conceptlabel, anchor=south] at (-1.4,1.25) {code specifies each step};
\node[conceptlabel, anchor=south] at (7.4,1.25) {model sets each step};
% the two blocks
\node[conceptbox, fill=blue!15,  minimum width=3cm] (wf) at (0,0) {Workflow};
\node[conceptbox, fill=green!15, minimum width=3cm] (ag) at (6,0) {Agent};
% distinctions, listed below each block (outside the node)
\node[conceptlabel, text width=4.4cm, anchor=north] at (0,-0.75)
  {steps fixed in code\\control flow known in advance\\predictable and cheaper\\chaining, routing, parallelization};
\node[conceptlabel, text width=4.4cm, anchor=north] at (6,-0.75)
  {model chooses next step at run time\\loops until task is completed\\flexible and open-ended\\higher cost, more room for error};
\end{tikzpicture}

\caption{The workflow--agent spectrum (Anthropic, 2024).
A \emph{workflow} routes the model through predetermined code paths; an \emph{agent} lets the model direct its own loop and tool use.
The added flexibility of an agent is bought with higher cost and more room for compounding error, so the durable advice is to reach for a workflow when the steps are known in advance and reserve full agency for open-ended tasks that genuinely need it.}
\label{figure:WorkflowAgentSpectrum}
\end{center}
\end{figure}

\section{The Capability Ladder}
\label{section:CapabilityLadder}

The course builds capability in a fixed order because each rung requires the one below it.\index{Capability ladder}
A researcher cannot \emph{build} an agent they cannot \emph{drive}, cannot \emph{orchestrate} what they cannot \emph{build}, and cannot \emph{evolve} what they cannot \emph{evaluate}.
The ladder has four rungs.

\begin{figure}[htb!]
\begin{center}
%% capability-ladder.tex — the four-rung capability ladder (Chapter 2). House conceptual style.
\input{figures/concept-style}
\begin{tikzpicture}[
    rung/.style={conceptbox, minimum width=3cm, font=\small\bfseries},
    lbl/.style={font=\scriptsize, align=left, text width=6.2cm},
]
\node[rung, fill=gray!15]   (op)  at (0,0)    {Operate};
\node[rung, fill=blue!15]   (bld) at (0,1.25) {Build};
\node[rung, fill=orange!15] (orc) at (0,2.5)  {Orchestrate};
\node[rung, fill=green!15]  (evl) at (0,3.75) {Evolve};
\draw[conceptflow] (op)  -- (bld);
\draw[conceptflow] (bld) -- (orc);
\draw[conceptflow] (orc) -- (evl);
\node[lbl, right=0.6cm of op]  {Terminal, shell, version control, model API};
\node[lbl, right=0.6cm of bld] {Tool definitions, the agent loop, an evaluation harness};
\node[lbl, right=0.6cm of orc] {Multi-agent workflows, retrieval, reliability};
\node[lbl, right=0.6cm of evl] {generate $\rightarrow$ evaluate $\rightarrow$ select $\rightarrow$ mutate};
\end{tikzpicture}

\caption{The capability ladder of the course.
Each rung is a prerequisite for the one above it.
The arrows indicate this dependency order.}
\label{figure:CapabilityLadder}
\end{center}
\end{figure}

\paragraph{Operate.}
The foundation is being able to drive an existing agent reliably and reproducibly.
The habit most worth instilling first is deceptively plain, stating the objective before the agent acts, which appears to track downstream proficiency more closely than knowledge of any particular feature (Anthropic, 2025).
This rung trains goal clarity before any tooling.
This means understanding the terminal, using version control, accessing a model through an API, and giving the agent instructions that produce predictable, revisable results.
Each of those has a purpose worth stating plainly before Chapter~\ref{chapter:AgenticRig} teaches the mechanics: the terminal is where an agent's actions actually run, version control records periodic snapshots of a project so that any edit an agent makes can be inspected and undone, and API access is what lets a model be called from a script rather than only from a chat window.
Many researchers skip this rung and never develop the habit of \emph{assessing modifications} the agent makes, a deficit that can grow more costly with each rung climbed.

There is a genuine tailwind here: as agentic tooling improves, operating the rig gets steadily easier in practice, so the mechanical barrier to entry keeps falling.
Asking Claude or Codex how to commit your work to git, once a real hurdle, is now trivial.
What does not get easier is the judgment underneath: understanding \emph{why} version control is of paramount importance, and when a commit is the right one, remains the researcher's to hold.
The same pattern recurs one rung up: getting an agent to write a piece of code for a task is increasingly dependable, but ensuring the build and architecture are conducive to discovery, rather than merely running, stays a cognitive challenge that the tooling does not carry for you.
And the escalation has a summit.
Beyond the rig and beyond the code sits the judgment no rung of tooling reaches: which questions are worth pursuing, what would count as a meaningful contribution, and what the results imply, for a field and beyond it.
That judgment stays the researcher's to hold; the promise of everything that follows is to clear the mechanical work away from it, not to make it for you.

\paragraph{Build.}
Building an agentic evolutionary framework  means implementing the loop and the tools from scratch rather than simply configuring a pre-existing system.
The user writes tool definitions, assembles a control loop, and decides how structured outputs will be formatted and validated.
Formatting and validation are less abstract than the words suggest: they amount to deciding that a proposal must come back in a fixed layout, say a candidate expression together with its parameter count, and then writing the check that rejects anything that does not parse.
An evaluation harness is also introduced at this rung: the agent must be judged by something more reliable than intuition.

\paragraph{Orchestrate.}
This rung involves composing multiple agents or steps into a reliable workflow.
It often entails failure handling, logging, memory, and the \emph{proposer–critic} pattern, in which one agent generates candidates and a second verifies or scores them.
The first two items in that list carry more weight than they look: a run that touches a few hundred candidates will meet a timed-out request, a malformed response, and a tool that silently returns nothing, and what separates an orchestration from a script is that it is written to expect these, record them, and carry on.
A reliable orchestration that can run unattended for hours is a prerequisite for the next rung: an evolutionary search cannot be trusted if it fails silently.

\paragraph{Evolve.}
Evolution is orchestration plus a mutation operator and a fitness function.
A system that generates candidates, evaluates them, and selects the best is an evolutionary algorithm whether or not it is called that.
The insight this rung turns on is that the evaluation harness introduced at the \emph{Build} rung, and carried through every rung since, \emph{is} the fitness function of the evolutionary search.
The practitioner has been building toward evolution without necessarily knowing it.
It is worth noting where the randomness that drives an evolutionary search naturally comes from.
A language model's output is a \emph{probability distribution} over next tokens, and sampling from it is, at its core, a mechanism for exploration: two runs of the same prompt can take different paths.
Modern evolutionary frameworks harness randomness through additional mechanisms (which candidates survive, how they are mutated and recombined), but it is worth noting that stochastic, divergent paths are native to this substrate rather than bolted on.

\section{The Two Spines}
\label{section:TwoSpines}

Two threads connect every rung of the ladder.\index{Two spines}
They are called \emph{spines} because every unit hangs from them.

\subsection{Spine 1: Evaluation Becomes Fitness}

You cannot trust an agent you cannot score, and you cannot evolve an artifact you cannot score automatically.
Because the model generates by plausible continuation rather than lookup, its output is fluent whether or not it is correct, which is precisely why trust must come from an external score rather than from the confidence of the text itself.
Evaluation is therefore introduced when a researcher first builds an agent and reused in every subsequent unit.
The key moment comes at the capstone pivot: the evaluation harness, which has been used to \emph{compare methods}, is now used to \emph{drive search}.
Nothing about the harness changes; its role changes.
One caveat travels with this promotion: a fixed evaluation loses discriminating power as capability rises (a capability-relative echo of Goodhart's law that returns as benchmark saturation in Chapter~\ref{chapter:Evaluation}), so the fitness function is best treated as an evolving artifact to be refreshed, not a fixed target.
This transition, evaluation becoming fitness, is the conceptual hinge of the course.
There is also a pedagogical reason evaluation is promoted to a spine rather than left to habit.
Not every competence is learned the same way: some (shaping context, reusing a project setup) tend to accrue with ordinary exposure, whereas the discernment to catch a wrong assumption or a plausible-but-false claim shows little tendency to grow on its own, an asymmetry consistent with Anthropic's analysis of how researchers use these tools (Anthropic, 2025).
Because discernment does not come free with practice, it must be taught deliberately and revisited at every rung.

\begin{principle}
\label{principle:EvalFitness}
A reusable validation harness built for comparison can serve as the fitness function for evolutionary search, but not every harness can, and not without care.
It must be \emph{automatic} (it runs without human judgment), \emph{leakage-resistant} (it cannot be gamed by artifacts that overfit the split it scores), cheap and stable enough to be queried many times over, and distinct from any final-test split held back for a one-shot audit.
A benchmark that is too costly, too noisy, non-stationary, unsafe under repeated use, or reserved for final reporting is not a drop-in fitness function.
\end{principle}

\begin{example}
Suppose an economist is evolving a forecasting rule for regional unemployment.
Dividing the historical record into three portions makes each criterion concrete: rules are fitted on a \emph{training} portion, the search proposes and scores them thousands of times against a \emph{validation} portion, and a \emph{test} portion covering the most recent years is locked away and consulted once, at the end, to report performance honestly.
\emph{Leakage} is what happens when that boundary is not airtight, for instance when a predictor is standardized using the whole series and so quietly carries information from the test years, and the rule then scores well for a reason that will not survive contact with the future.
Cost matters just as sharply: an evaluation taking ten minutes is perfectly workable for comparing three methods by hand, and completely infeasible for a search that must score thousands of candidates.
\end{example}

How that automatic score is computed depends on the task, along a spectrum of \emph{evaluability}.
At one end are \emph{verifiable} problems, where an unambiguous oracle settles quality: the density of a sphere packing, whether a program passes its tests against a defined interface, the error of a fitted model.
At the other are problems of \emph{judgment}, where no formula exists and a capable model must act as the evaluator, a \emph{critic} scoring an artifact against a rubric: the soundness of a proof at the research frontier, the quality of a written argument, the promise of an idea.
Which end a task falls on is a property of the problem, not a preference.
The verifiable end is where evaluation-becomes-fitness is cleanest, and where automated search has so far reached furthest; the judgment end is softer and more easily gamed, yet as models grow more capable it becomes usable, so strict formula-based evaluation is no longer the only norm.
Chapter~\ref{chapter:Evaluation} develops both ends.

\subsection{Spine 2: Context Engineering}

An LLM is a function of its context.
What the model proposes, generates, or decides depends entirely on what it reads.
Context engineering, the deliberate construction of what goes into the context window, is therefore the primary lever available to the agent builder.

This lever has itself matured through three levels.
\emph{Prompting}, the mode of early and web-based use around 2023, was about which question to ask and how to phrase it.
\emph{Context engineering} arrived with large windows and persistent profiles: the question became how to fill a hundred-thousand- or million-token window well, where to position each element within it, when caching earlier tokens pays off, and when to clear the window rather than let it accumulate.
\emph{Agentic context engineering} is the current shift, and the reason this spine is named with that adjective in mind: the context is increasingly assembled not by the researcher but by an agent, such as Claude or Codex, working on the researcher's own machine, which is what motivates giving that agent access to one's files, tools, and knowledge.
The questions move accordingly, toward how to keep a file hierarchy an agent can navigate, how to make a corpus searchable with ordinary unix tools (\code{grep}, \code{sed}, \code{awk}), which file formats are cheap in tokens and time, and how to curate the tools, instruction files (\code{CLAUDE.md}, \code{AGENTS.md}), skills, and connectors an agent draws on.
Those unix tools are decades old and deliberately unglamorous: \code{grep} reports every line matching a pattern across a whole directory tree, so an agent able to run it can search a corpus of a thousand files while pulling only the matching lines into its window.
The instruction files are equally plain, ordinary text kept at a project's root and read at the start of a session, so that standing conventions need not be restated each time.
Chapters~\ref{chapter:AgenticRig} and~\ref{chapter:ContextEngineering} take up these subtleties; here it is enough to feel the richness, that the lever has moved from wording a prompt to shaping the environment an agent builds its own context from.

This thread begins with prompt engineering and agentic context management, deepens into retrieval-augmented grounding in Chapter~\ref{chapter:Retrieval}, and reaches its culmination in the evolutionary capstone: the LLM mutation operator is most effective when its context contains the lineage of what has been tried, the domain priors, and the current population's best candidates.
Evolution is context engineering in a loop.

\section{Why Now}

Agentic AI systems have existed in various forms for decades: rule-based expert systems, reinforcement learning agents, robotic planning systems.
What makes the current moment distinct is the arrival of \emph{general-purpose tool-using language models}: systems that can read natural language instructions for a new task, decide which tools to invoke, interpret unstructured results, and adapt mid-task without domain-specific engineering.
Three capabilities combine to make this shift significant.

\paragraph{Instruction following.}
Modern language models follow complex, nuanced instructions expressed in natural language with high reliability.
An agent builder need not program the model's behavior in code; they write a system prompt.
This lowers the barrier to specifying what an agent should do: a matter of clear description, not programming.
A geologist who wants an agent to download seismic records, filter them by magnitude, and produce a summary figure writes that goal in plain English, not in a programming language the geologist may not speak fluently.

\paragraph{Tool use.}
Tool-calling capability, now standard in frontier models, allows a language model to take actions in the world (reading files, running code, searching the web) and incorporate the results into subsequent reasoning.
The model can, in principle, perform the same actions a human researcher performs at a computer.

\paragraph{Structured output.}
Models can be constrained to produce output in a specified schema (JSON, typed records, code blocks), allowing the output of one agent to serve as the structured input of the next.
JSON is simply a text format for labelled values, and the constraint buys something mechanical: a downstream step can read a field by name instead of hunting for a number inside a paragraph of prose.
This is the glue of multi-agent orchestration.

Together, these capabilities make a language model a plausible substrate for a research assistant that can be directed, composed, and evaluated programmatically.
The course is concerned with doing this \emph{well}: reliably, reproducibly, and honestly.
Doing it well requires a working rig (the terminal, version control, and API access on which everything else rests), and that is where the next chapter begins.

\begin{remark}
A fair question hangs over everything that follows: if the platform keeps improving (automatic context management, better tool use, longer windows, models that reason before they act), does the researcher's own engineering simply get absorbed?
Part of it does, and should: generic prompt-crafting and transport-level tuning (how context is cached and ordered, when to compact a conversation) are being standardized into the platform, and a wise builder lets them be.
What does not compress away is the part specific to your problem: the design of an evaluation that measures what \emph{you} mean by better, the curation of a corpus that encodes your field's judgment, and the architecture of the loop that composes these into something trustworthy.
This is why the chapters ahead spend their effort on evaluation, grounding, and orchestration rather than on prompt wording: the durable skill is building the parts a general system cannot know to build for you.
\end{remark}

\section*{Further Reading}

\emph{The engineering-blog and product references below are Fall~2026 snapshots of a fast-moving surface; titles and locations may have since changed.}

\begin{small}
\begin{enumerate}
\item Yao, S., et al., ``ReAct: Synergizing Reasoning and Acting in Language Models,'' \emph{ICLR}, 2023 (arXiv:2210.03629) — the canonical formalization of the observe–reason–act loop.
\item Anthropic, \emph{Claude's Tool Use Guide}, available at \href{https://docs.anthropic.com}{docs.anthropic.com} — the practical tool-calling interface used throughout the course.
\item Wang, L., et al., ``A Survey on Large Language Model Based Autonomous Agents,'' \emph{Frontiers of Computer Science}, 2024 (arXiv:2308.11432) — a comprehensive taxonomy of agent architectures.
\item Anthropic, \emph{AI Fluency} (the four-properties and the four-competency ``4Ds'' framing) together with its accompanying analysis of agent-usage patterns, 2025 — the source for this chapter's substrate-properties framing and for the observation that goal clarity and discernment track proficiency more than feature knowledge does.
\item Anthropic, ``Building Effective Agents,'' \emph{Anthropic Engineering}, 2024 — the workflow-versus-agent distinction adopted here, with a catalogue of composable workflow patterns and the counsel to prefer the simplest design a task allows.
\end{enumerate}
\end{small}
% EVOLVE-BLOCK-END
